\section{Conclusion}
\label{sec:conclusion}

This paper set out to close the step between a predictive signal and an executable
portfolio, and to determine which parts of a regime-conditional framework are
responsible for its behaviour. Both objectives are met, though not in the direction
originally anticipated.

The framework, evaluated on ten liquid multi-asset instruments over sixteen years under
a walk-forward protocol with transaction costs, attains a Sharpe ratio of 0.738 against
0.532 for an equal-weighted benchmark and reduces maximum drawdown from $-36.1$ to
$-19.0$ per cent, while earning a lower compound return. It outperforms in all four
major stress episodes in the sample, improves on drawdown in 15 of 16 calendar years,
and beats the benchmark on both measures in all 18 specifications examined. The Sharpe
improvement is consistent but not statistically significant: its bootstrap interval
includes zero.

Ablation locates the source of that behaviour precisely. At identical average exposure,
timing risk by the identified market state improves both the Sharpe ratio and the
maximum drawdown with no difference in mean return, which is the signature of a genuine
risk-timing effect. The form of the risk budget proves more consequential than its
calibration: expressed as a fraction of the portfolio's own volatility the layer is
active and effective, while the absolute volatility targets of the original
specification bind in one rebalance out of 192 and leave it inert.

The signal fusion layer, which gave the framework its name, does not contribute.
Removing it improves the Sharpe ratio, and the information coefficients of both
composites are significantly negative over the full sample. The regime-conditional
weighting schedule is worse than neutral: it was specified on the hypothesis that
market-state information gains importance under stress, and the data reject that
hypothesis in the opposite direction at $p = 0.008$. This null is reported in full,
with the coefficients that produce it, because a component-level negative result is more
useful than an aggregate number that conceals which parts of an architecture are doing
the work.

The study's contribution is therefore a specification and a method of testing it: a
two-step separation of composition from total risk that makes each layer independently
falsifiable, an ablation design that isolates risk timing from risk level by matching
average exposure, and a documented account of two failure modes---an inert absolute risk
budget, and a signal-to-allocation conversion in which mismatched units rather than
information determine the solution. The framework is useful within stated bounds: it
requires a universe with genuine dispersion in asset volatility, it trades return for
drawdown at a measured rate, and it should be evaluated against a volatility-matched
benchmark rather than an unlevered one. Outside those bounds, as the mega-capitalisation
results show, it does not add value.
